\documentclass[10pt,twocolumn,twoside,letterpaper]{article}

% ---------------------------------------------------------------- layout (workshop style, two columns)
\usepackage[letterpaper,top=0.9in,bottom=0.9in,left=0.75in,right=0.75in,columnsep=0.25in]{geometry}
\usepackage[T1]{fontenc}
\usepackage{times}
\usepackage[scaled=0.86]{helvet}
\usepackage{microtype}
\usepackage{amsmath,amssymb}
\usepackage{booktabs}
\usepackage{tabularx}
\usepackage{multirow}
\usepackage{xcolor}
\usepackage{graphicx}
\usepackage{enumitem}
\usepackage[font=small,labelfont=bf,skip=4pt]{caption}
\usepackage[numbers,sort&compress]{natbib}
\usepackage{tikz}
\usetikzlibrary{arrows.meta,positioning,fit}
\usepackage{fancyhdr}
\usepackage{titlesec}
\usepackage[hidelinks]{hyperref}

\titleformat{\section}{\normalfont\large\bfseries}{\thesection}{0.6em}{}
\titleformat{\subsection}{\normalfont\normalsize\bfseries}{\thesubsection}{0.6em}{}
\titlespacing*{\section}{0pt}{1.0ex plus 0.4ex}{0.6ex}
\titlespacing*{\subsection}{0pt}{0.8ex plus 0.3ex}{0.4ex}
\setlist{nosep,leftmargin=*}
\setlength{\parskip}{0pt}
\setlength{\bibsep}{1pt}
\renewcommand{\bibfont}{\footnotesize}

\pagestyle{fancy}
\fancyhf{}
\fancyhead[LE,RO]{\footnotesize\thepage}
\fancyhead[RE,LO]{\footnotesize\textit{Onus: claim-to-evidence maps for research papers}}
\renewcommand{\headrulewidth}{0.4pt}

% placeholders for numbers that are not in yet (held-out test runs)
\newcommand{\TBD}[1]{\textcolor{red}{\textbf{[#1]}}}
\newcommand{\onus}{\textsc{Onus}}
\newcommand{\verdict}[1]{\texttt{#1}}

\title{\vspace{-2.2em}\Large\bfseries \onus{}: Claim-to-Evidence Maps for Research Papers,\\with a Measured Error Rate}
\author{Avinash Ranganath\\[1pt]{\small \href{https://avinashranganath.com}{avinashranganath.com} \quad AI4MetaScience Hackathon 2026, Task 3 (claim-to-evidence linker)}}
\date{\vspace{-2.5em}}

\begin{document}
\maketitle
\thispagestyle{fancy}

% ================================================================ abstract
\begin{abstract}
\noindent
Reviewers rarely trace each claim in a paper to the evidence behind it, so confident wording passes for
support. We present \onus{}, which turns a paper (PDF, arXiv or OpenReview identifier) into a \emph{claim
map}: every major claim is quoted verbatim with its location, linked to the table cells, passages, checklist
answers and code lines that bear on it, and given one of five verdicts (\verdict{supported},
\verdict{overstated}, \verdict{untraced}, \verdict{contradicted}, \verdict{not\_checkable}) with a reason.
Language models only \emph{point}; arithmetic, table comparisons, checklist rules and verdict precedence are
code. Where code is available, \onus{} matches the paper's implementation facts to code lines and re-runs the
code in a sandbox. We evaluate on the 247 Agents4Science 2025 submissions with a frozen, fingerprinted
version and held-out test sets run once. On held-out papers, \onus{} flags half as many claims as a
single-prompt baseline (15\% vs.\ 29--32\% of claims on unmutated papers) at statistically similar recall of
planted errors: 70\% [58--80] vs.\ 76\% [64--85] on test\_1 and 75\% [62--84] vs.\ 78\% [66--87] on the
replication set test\_2. Its flags are less precise on blind labels (67\% vs.\ 89\%), mostly because of
\verdict{untraced} verdicts. A deployed web app analyses a 15-page paper in 5--8 minutes for about
\texteuro0.10.
\end{abstract}

% ================================================================ 1 introduction
\section{Introduction}
Peer review assumes that a paper's claims are backed by its own results. In practice, wording often outruns
evidence: a method ``consistently outperforms'' baselines on two of five datasets, an abstract number differs
from its table, or promised code is missing \citep{lipton2019troubling}. The problem grows as language models
make fluent papers cheap. At Agents4Science 2025, where AI agents were primary authors \emph{and} reviewers
\citep{bianchi2025agents4science}, 42 of 247 submissions point to code that is not there, and the
conference's audit of its 48 accepted papers found 6 critical paper--code mismatches.

\onus{} produces a claim map that a reviewer can inspect and dispute, not a score to trust. Our
contributions: (i) a claim map that joins text, numbers, checklist and code, with located evidence, the
checks that ran, a verdict with a reason, stability across runs and a ledger of what was \emph{not} checked;
(ii) a ``code computes, the model points'' design in which every model quote, evidence identifier and cited
code line is verified; (iii) an evaluation that could not be tuned on its test data, reporting recall,
precision, a baseline and stability as they came out; and (iv) a deployed system for arXiv, OpenReview or
uploaded papers, with live code checks.

\textbf{Related work.} Rule-based tools recompute one error type at scale: statcheck for $p$-values
\citep{nuijten2016prevalence}, GRIM for impossible means \citep{brown2017grim}. Reproducibility work produced
checklists and code policies \citep{pineau2021improving}, benchmarks for reproducing results from code or
papers \citep{siegel2024corebench,starace2025paperbench}, and repository-level audits of agent-written papers
\citep{shen2026beyond}. FEVER and SciFact frame verification as retrieving and labelling evidence
\citep{thorne2018fever,wadden2020fact}; LLM reviewers overlap with human reviews \citep{liang2024can}, and
automated pipelines now write and review papers \citep{lu2024aiscientist}. \onus{} recomputes rather than
judges, works per claim across error types, and uses the paper's own tables, checklist and code as the
evidence base.

% ================================================================ 2 system
\section{The \onus{} System}
\label{sec:system}

\begin{figure*}[t]
\centering
\begin{tikzpicture}[
  node distance=0.30cm,
  box/.style={draw, rounded corners=2pt, align=center, font=\footnotesize, minimum height=1.0cm,
              text width=1.95cm, fill=#1},
  box/.default={gray!8},
  arr/.style={-{Stealth[length=4pt]}, thick}]
\node[box=blue!6]   (in)  {\textbf{Paper}\\PDF / arXiv /\\OpenReview\\+ code};
\node[box, right=of in]  (ing) {\textbf{Ingest}\\text, sections,\\tables, captions,\\checklist};
\node[box, right=of ing] (cl)  {\textbf{Claims}\\LLM extracts;\\quotes verified;\\typed, located};
\node[box, right=of cl]  (ev)  {\textbf{Evidence}\\explicit refs +\\BM25 $\cup$ embed.;\\LLM selects};
\node[box, right=of ev]  (ck)  {\textbf{Checks}\\numbers, dominance,\\strength, scope,\\checklist, code};
\node[box, right=of ck]  (vd)  {\textbf{Verdicts}\\precedence rules;\\adjudicator may\\only downgrade};
\node[box=green!7, right=of vd]  (out) {\textbf{Claim map}\\3 runs $\to$\\stability +\\coverage ledger};
\foreach \a/\b in {in/ing,ing/cl,cl/ev,ev/ck,ck/vd,vd/out} \draw[arr] (\a) -- (\b);
\node[draw, dashed, rounded corners=2pt, font=\footnotesize, align=center, text width=4.4cm,
      below=0.3cm of ck, xshift=-1.2cm] (code)
      {\textbf{Code checks (code-v1)}: implementation facts vs.\ code lines;
       re-run in an isolated sandbox};
\draw[arr, dashed] (code.north) -- ([xshift=-1.2cm]ck.south);
\end{tikzpicture}
\caption{The \onus{} pipeline. Solid boxes are the frozen pipeline (v1.0.0); the dashed box is the separately
versioned code-check stage. Model calls are limited to extraction, selection, a few classifications and
adjudication; all comparisons and verdict rules are code.}
\label{fig:pipeline}
\end{figure*}

\textbf{Claims and evidence.} The PDF is parsed into sections, sentences, tables (as cells), captions and the
venue checklist. An extraction model lists claims from the abstract, introduction and conclusion as verbatim
quotes; a quote not found in the text is dropped. Every checklist answer becomes a claim without a model call.
Explicit references (``Table~3'') are resolved first, and an unresolved one is recorded as \emph{dangling};
BM25 \citep{robertson2009bm25} and an embedding model \citep{wang2024multilingual} add candidates, and a model
labels each as supporting, partial or refuting, citing only identifiers it was shown.

\textbf{Checks and verdicts.} \emph{Numbers}: the model maps each number in a claim to table cells and an
operation; code computes and compares within a precision-aware tolerance. \emph{Dominance}: code checks that
``ours'' beats every listed baseline in every listed cell. \emph{Strength} and \emph{scope}: rules require
evidence proportional to wording (``significantly'' needs a test; ``robust'' needs two or more settings) and
check strong claims against stated limitations. \emph{Checklist audit}: each ``Yes'' is tested against the
paper. A precedence table maps outcomes to verdicts; a reasoning model then reviews every flag but may only
confirm or downgrade it, citing evidence. In batch mode the pipeline runs three times; the majority verdict
wins and agreement is reported as \emph{stability}.

\textbf{Code checks.} Up to eight implementation facts (datasets, algorithms, metrics, hyperparameters) are
extracted with verified quotes and rated present, absent or contradicted against the code, citing
\texttt{path:line} references that are checked against the files. An agent re-runs the code in a sandbox
(network on for installation with no secrets, off for the run), capped in steps, minutes and cost; reported
values are compared with the paper's numbers by code.

\textbf{Models and serving.} An extraction model (GLM-5.3-Flash), a judge (Kimi-K3.1-Pro), a fast classifier
(Gemma-4-31B) and multilingual-e5-large-instruct embeddings. The app (FastAPI, React) uses Berget AI, an EU
provider that does not train on inputs; its report places the PDF beside the claim list and highlights each
quote and evidence item in the paper. Held-out batches used the same models through a second provider because
of credit limits; on dev papers both paths gave the same flagged/not-flagged decision on 34 of 34 matched
claims.

\begin{table*}[t]
\centering\footnotesize
\setlength{\tabcolsep}{4pt}
\resizebox{\textwidth}{!}{\begin{tabular}{@{}lcccccc@{}}
\toprule
 & \multicolumn{2}{c}{eval (pre-freeze, run \#2)} & \multicolumn{2}{c}{test\_1 (frozen, once)} &
   \multicolumn{2}{c}{test\_2 (frozen, once)} \\
\cmidrule(lr){2-3}\cmidrule(lr){4-5}\cmidrule(lr){6-7}
Measure & \onus{} & Baseline & \onus{} & Baseline & \onus{} & Baseline \\
\midrule
\textbf{Flag rate, unmutated papers} & \textbf{21\%} [17--25] & 31\% [26--35] & \textbf{15\%} [12--18] & 29\% [25--33] & \textbf{15\%} [13--18] & 32\% [29--36] \\
Recall, planted errors & 25/30 (83\%) [66--93] & 24/30 (80\%) [63--90] &
  44/63 (70\%) [58--80] & 48/63 (76\%) [64--85] & 44/59 (75\%) [62--84] & 46/59 (78\%) [66--87] \\
Collateral flag rate        & 2\% [1--3]   & 4\% [3--6]   & 2\% [1--3] & 6\% [5--8] & 3\% [2--4] & 4\% [3--6] \\
Precision (labelled flags)  & ---          & ---          & 39/58 (67\%) [54--78] & 24/27 (89\%) [72--96] & --- & --- \\
Stability (3 runs) & \multicolumn{2}{c}{1 run} & \multicolumn{2}{c}{84\% [81--87]} &
  \multicolumn{2}{c}{84\% [81--87]} \\
\bottomrule
\end{tabular}}
\caption{Flag rates, planted-error recall and precision (Wilson 95\% intervals). Eval: 20 papers, 1 run;
test: 30 papers each, frozen version, 3 runs. Test numbers are reported as they came out.}
\label{tab:results}
\end{table*}

% ================================================================ 3 evaluation
\section{Evaluation Protocol}
\label{sec:protocol}

The 247 Agents4Science 2025 submissions (48 accepted) were split once with a fixed seed, before any pipeline
code existed, into dev (10), eval (20), test\_1 (30), test\_2 (30) and a remainder. Development and tuning used
only dev; held-out splits stayed in an encrypted archive outside the repository. After two eval runs (of at
most three) the version was frozen with a fingerprint over models, prompts, rubrics and pipeline source, and
each test split was run once; every held-out run was logged publicly before it started.

\textbf{Measures.} (1)~\emph{Recall on planted errors} \citep{demillo1978hints}: five mutation types are
planted where detectable: a changed abstract number (caught if \verdict{contradicted}), a checklist answer
flipped to ``Yes'', a deleted table, a specific claim widened to a universal one, and deleted error bars.
The \emph{collateral flag rate} counts new flags away from the mutation. (2)~\emph{A single-prompt baseline}
judges the same claims with the same judge model; since recall flatters a system that flags a lot, each
system's flag rate on unmutated papers is reported beside it. (3)~\emph{Precision}: a blind, shuffled sample
of 60 \onus{} and 30 baseline flags from test\_1, labelled correct or incorrect. Labels were drafted by an AI
assistant from the paper text, blind to the system, with a reason each; the author reviewed all 90 and
changed none. (4)~\emph{Stability} across three runs. Intervals are Wilson 95\% \citep{wilson1927probable}.

% ================================================================ 4 results
\section{Results}


\textbf{Main result.} On every split, eval included, \onus{} flags roughly half as many claims as the baseline
at recall within the baseline's interval (Table~\ref{tab:results}): 44 of 63 planted errors vs.\ 48 of 63 on
test\_1, and 44 of 59 vs.\ 46 of 59 on test\_2, with a third of the baseline's collateral flags on test\_1. We
do not claim a recall difference at this sample size; the baseline reaches its recall by flagging about twice
as many claims. 84\% of claims had the same verdict in all three runs.

\textbf{Where it falls short.} Precision does not favour \onus{}: 39 of 58 flags are correct against 24 of 27
for the baseline. By verdict, \verdict{overstated} is right 19 of 23 times, \verdict{contradicted} 2 of 4 and
\verdict{untraced} only 18 of 31, the main source of wrong flags. By mutation type, \onus{} matches or beats
the baseline on deleted error bars, flipped checklist answers and widened scope, but catches fewer changed
abstract numbers on test\_1 (4/11 vs.\ 10/11); deleted tables are rarely caught by either system (2 of 9 and 1
of 9). \onus{} is the more conservative system, not the more precise one.

\textbf{Code checks and corpus findings.} On dev, the static check agreed with the conference's code audit on
all three labelled papers with a verdict; in one, a claim about a real patient cohort was contradicted by the
code's own docstring (``Synthetic dataset generation''). Across 90 analysed papers, 68 have code; re-runs were
attempted on 48 and 3 reproduced the paper's values. Checklist ``Yes'' answers are often unbacked: 12 of 18 for
theory, 35 of 69 for broader impacts, 18 of 56 for compute, 19 of 81 for limitations. A paper's flagged share is
unrelated to its review scores (Spearman $\rho=-0.05$ AI, $0.10$ human) and to acceptance (difference 0.9
points, $p=0.82$).

\textbf{Cost and examples.} An analysis costs about \texteuro0.10 per run (15--30 claims per paper). On
\emph{Attention Is All You Need}, \onus{} marked the claim of a new state of the art ``after being trained for
as little as twelve hours'' as \verdict{overstated}: the twelve-hour figure belongs to the base model, the
state-of-the-art score to the larger one. On a 2026 arXiv paper, it found the repository through a link broken
across lines and reported that the code covers the simulations but none of the three real-data case studies.

% ================================================================ 5 limitations + conclusion
\section{Limitations and Conclusion}
About 60 planted errors per test split give recall intervals of about $\pm$11 points. Planted errors are
text-level and cover five types. Precision labels were AI-drafted and reviewed by one author, without
inter-rater agreement. Accepted papers are camera-ready versions. Re-runs are CPU-only and time-capped, and
live analyses run once.

\onus{} shows that a claim-level map joining text, tables, checklists and code can be produced cheaply, and
that its error rate can be measured honestly. It flags half as many claims as a single-prompt baseline at
similar recall, but its flags are less precise, mostly because of \verdict{untraced}, which is the first
thing to improve. It is meant to route a reviewer's attention, not replace the reviewer. Code, protocol and
run log: \url{https://github.com/nash911/onus}; live app: \url{https://www.onuscheck.com}.

{\small
\bibliographystyle{plainnat}
\bibliography{refs}
}

\end{document}
